\documentclass[mylettercfg2.tex,11pt]{ucsdletter}% to use a config file other than mylettercfg.tex
%you need to include the filename in the options, like
%\documentclass[myotherlettercfg.tex,11pt]{ucsdletter}
\usepackage{amsmath}

\begin{document}

\date{August 11, 2024}

\begin{letter}{Dear Editors:}

\opening{}

On behalf of my coauthor Ethan N. Epperly and myself, I am writing to submit a revised version of our manuscript ``Randomized Kaczmarz with tail averaging'' for consideration in Applied Computational and Harmonic Analysis.

% The manuscript offers a new analysis and extension of the randomized Kaczmarz (RK) method.
% RK is a simple stochastic algorithm that converges exponentially fast to the solution of a consistent least-squares problem $\min_{\mathbf{x}} \lVert \mathbf{A} \mathbf{x} - \mathbf{b} \rVert^2$. RK is an important method to understand because it is scalable --- it accesses just a small fraction of the data in each iteration. Moreover, the analysis of RK has already facilitated an improved understanding of the capabilities and limitations of stochastic gradient descent.

% Despite the apparent simplicity of RK, for a long time it has been unclear how to extend the method to handle inconsistent least-squares problems using only row-wise access to the data. The method as originally presented does not converge in the inconsistent case, and there have been several attempts to improve the convergence behavior, none of them wholly satisfactory. Here, we provide a simple solution for inconsistent systems based on tail-averaging, and we support this recommendation with numerical experiments, theoretical error bounds, and a detailed discussion of the theoretical optimality of our proposed tail-averaged randomized Kaczmarz (TARK) method. 

% After emailing our preprint to several RK experts, they have indicated the idea of tail-averaging is both natural and computationally efficient. We anticipate that it will be a serious contribution to the literature that will be appreciated by the readers of ACHA.

Thank you for your consideration of the manuscript.

\closing{Sincerely,}

\end{letter}

\end{document}
